\chapter{Appendix}
\label{ch:appendix}

\section{Hyperparameter Search Spaces}
\label{sec:app-grids}

Table~\ref{tab:grids} lists the parameters searched for each model in the nested
grid search described in Section~\ref{sec:cv}. The three grids contain $12$, $36$
and $12$ combinations respectively, and since every combination is evaluated over
the full $30$ cross-validation fits, the search amounts to $1{,}800$ model fits in
total. The winning configurations are reported in Table~\ref{tab:selected} of
Section~\ref{sec:res-performance}.

\begin{table}[htbp]
  \centering
  \caption{Hyperparameter search spaces per model.}
  \label{tab:grids}
  \begin{tabular}{lll}
    \toprule
    \textbf{Model} & \textbf{Parameter} & \textbf{Values searched} \\
    \midrule
    \multirow{2}{*}{Logistic Regression}
      & \texttt{penalty}            & $\ell_1$, $\ell_2$ \\
      & \texttt{C}                  & $0.001$, $0.01$, $0.1$, $1.0$, $10.0$, $100.0$ \\
    \midrule
    \multirow{4}{*}{Random Forest}
      & \texttt{n\_estimators}      & $300$ (fixed) \\
      & \texttt{max\_depth}         & $3$, $5$, $10$ \\
      & \texttt{max\_features}      & \texttt{sqrt}, \texttt{log2}, $0.3$, $0.2$ \\
      & \texttt{min\_samples\_leaf} & $3$, $5$, $10$ \\
    \midrule
    \multirow{5}{*}{Gradient Boosting}
      & \texttt{n\_estimators}      & $200$ (fixed) \\
      & \texttt{learning\_rate}     & $0.05$, $0.1$ \\
      & \texttt{max\_depth}         & $1$, $2$, $3$ \\
      & \texttt{min\_samples\_leaf} & $1$, $5$ \\
      & \texttt{subsample}          & $0.8$ (fixed) \\
    \bottomrule
  \end{tabular}
  \floatfoot{Source: own.}
\end{table}

\section{Native Importance Rankings}
\label{sec:app-native}

Tables~\ref{tab:native-lr} to~\ref{tab:native-gb} list the top $15$ features by
native importance for each model, complementing the leading entries reported in
Section~\ref{sec:res-fi}. For the logistic regression the values are signed
coefficients, where a positive sign pushes the prediction toward bankruptcy. For
the two ensembles they are Gini impurity importances, which sum to one across all
features.

\begin{table}[htbp]
  \centering
  \caption{Native importance, logistic regression: top 15 signed coefficients.}
  \label{tab:native-lr}
  \small
  \begin{tabular}{rlr}
    \toprule
    & \textbf{Feature} & \textbf{Coefficient} \\
    \midrule
    1  & Debt ratio \%                          & $+14.493$ \\
    2  & ROA(B) before interest and depreciation after tax & $-9.558$ \\
    3  & ROA(A) before interest and \% after tax & $-5.193$ \\
    4  & Cash/Total Assets                      & $-4.738$ \\
    5  & Retained Earnings to Total Assets      & $+4.053$ \\
    6  & Total Asset Turnover                   & $-2.758$ \\
    7  & Total assets to GNP price              & $-2.678$ \\
    8  & Interest-bearing debt interest rate    & $-1.624$ \\
    9  & ROA(C) before interest and depreciation before interest & $-1.314$ \\
    10 & Tax rate (A)                           & $-1.219$ \\
    11 & After-tax net Interest Rate            & $+1.192$ \\
    12 & Liability-Assets Flag                  & $-0.997$ \\
    13 & Cash Turnover Rate                     & $-0.764$ \\
    14 & Fixed Assets Turnover Frequency        & $+0.569$ \\
    15 & Operating Profit Rate                  & $+0.565$ \\
    \bottomrule
  \end{tabular}
  \floatfoot{Source: own.}
\end{table}

\begin{table}[htbp]
  \centering
  \caption{Native importance, random forest: top 15 Gini impurity importances.}
  \label{tab:native-rf}
  \small
  \begin{tabular}{rlr}
    \toprule
    & \textbf{Feature} & \textbf{Importance} \\
    \midrule
    1  & Persistent EPS in the Last Four Seasons     & $0.092$ \\
    2  & Net Income to Stockholder's Equity          & $0.073$ \\
    3  & Net profit before tax/Paid-in capital       & $0.059$ \\
    4  & Per Share Net profit before tax (Yuan)      & $0.053$ \\
    5  & Continuous interest rate (after tax)        & $0.046$ \\
    6  & Borrowing dependency                        & $0.041$ \\
    7  & Net Income to Total Assets                  & $0.032$ \\
    8  & Total income/Total expense                  & $0.028$ \\
    9  & Non-industry income and expenditure/revenue & $0.027$ \\
    10 & Net worth/Assets                            & $0.025$ \\
    11 & Debt ratio \%                               & $0.025$ \\
    12 & Equity to Liability                         & $0.024$ \\
    13 & Interest Expense Ratio                      & $0.023$ \\
    14 & Degree of Financial Leverage (DFL)          & $0.022$ \\
    15 & Interest Coverage Ratio (Interest expense to EBIT) & $0.020$ \\
    \bottomrule
  \end{tabular}
  \floatfoot{Source: own.}
\end{table}

\begin{table}[htbp]
  \centering
  \caption{Native importance, gradient boosting: top 15 Gini impurity importances.}
  \label{tab:native-gb}
  \small
  \begin{tabular}{rlr}
    \toprule
    & \textbf{Feature} & \textbf{Importance} \\
    \midrule
    1  & Persistent EPS in the Last Four Seasons     & $0.247$ \\
    2  & Borrowing dependency                        & $0.105$ \\
    3  & Non-industry income and expenditure/revenue & $0.079$ \\
    4  & Net Income to Stockholder's Equity          & $0.052$ \\
    5  & Continuous interest rate (after tax)        & $0.048$ \\
    6  & ROA(C) before interest and depreciation before interest & $0.035$ \\
    7  & Degree of Financial Leverage (DFL)          & $0.034$ \\
    8  & Interest Coverage Ratio (Interest expense to EBIT) & $0.027$ \\
    9  & ROA(B) before interest and depreciation after tax & $0.026$ \\
    10 & Cash/Total Assets                           & $0.025$ \\
    11 & Inventory/Working Capital                   & $0.025$ \\
    12 & Equity to Liability                         & $0.023$ \\
    13 & Debt ratio \%                               & $0.021$ \\
    14 & Interest Expense Ratio                      & $0.021$ \\
    15 & Net worth/Assets                            & $0.015$ \\
    \bottomrule
  \end{tabular}
  \floatfoot{Source: own.}
\end{table}

\section{Complexity and Overfitting Curves}
\label{sec:app-complexity}

The following figures provide the overfitting check referenced in
Section~\ref{sec:cv}. For each model one complexity parameter is swept while all
remaining parameters are held at their grid-search optimum and training and
validation $F_2$ are traced across that sweep. The selected configurations of
Table~\ref{tab:selected} sit at or before the point where the validation curve turns
down or stagnates, which is what rules out overfitting at the chosen settings. As noted in
Section~\ref{subsec:gb}, the depth axis is not comparable across the random forest
and gradient boosting.

\begin{figure}[htbp]
  \centering
  \includegraphics[width=0.75\textwidth]{complexity_curve_LogisticRegression.png}
  \caption{Logistic regression: training vs.\ validation $F_2$ across the inverse
    regularisation strength $C$, swept from $0.01$ to $100$ on a logarithmic axis.}
  \label{fig:cc-lr}
  \floatfoot{Source: own.}
\end{figure}

\begin{figure}[htbp]
  \centering
  \includegraphics[width=0.75\textwidth]{complexity_curve_RandomForest.png}
  \caption{Random forest: training vs.\ validation $F_2$ across
    \texttt{max\_depth}, swept from $1$ to $20$.}
  \label{fig:cc-rf}
  \floatfoot{Source: own.}
\end{figure}

\begin{figure}[htbp]
  \centering
  \includegraphics[width=0.75\textwidth]{complexity_curve_GradientBoosting.png}
  \caption{Gradient boosting: training vs.\ validation $F_2$ across
    \texttt{max\_depth}, swept from $1$ to $8$.}
  \label{fig:cc-gb}
  \floatfoot{Source: own.}
\end{figure}
